% ── Section V-B: Fairness ─────────────────────────────────────────────────────
% Drop into the Evaluation section after the single-flow (T90) results.

\subsection{Fairness Under Competition}
\label{sec:fairness}

Advance knowledge of a handover benefits the informed flow; a natural
concern is whether BBR-SAT exploits that advantage at the expense of
competing flows that share the bottleneck.
We evaluate two scenarios using the shared-link simulator described in
Section~\ref{sec:methodology}:
\textbf{F1}~---~two flows both crossing a LEO$\to$GEO handover at $T=30$\,s,
and
\textbf{F3}~---~two flows that start simultaneously on a GEO beam
and reach steady state before BBR-SAT re-anchors its BDP context at
$T=30$\,s.
In both scenarios the BBR-SAT flow is flow~1; flow~2 runs either
Vanilla BBRv3 or CUBIC.
Per-second throughput and Jain's fairness index~\cite{jain1984} are
recorded for the full 90-second simulation; post-event averages
($t > 30$\,s) are reported in Table~\ref{tab:exp2}.
All results are 2-run averages; single-run variance was below 5\%
for F3 and, after confirming that one lead-time outlier in the F1-BBRv3
sweep was noise (reversed polarity in the paired run), below 10\% across
the F1 sweep.

\subsubsection{F1: shared LEO$\to$GEO handover}

Figure~\ref{fig:exp2_f1} shows the 90-second throughput trace for
BBR-SAT and BBRv3 when both flows cross the handover with zero lead time.
Before the handover ($t < 30$\,s) both flows share the 50\,Mbps LEO
beam equitably, each sustaining roughly 24\,Mbps.
At $t = 30$\,s the link drops to 10\,Mbps and the RTT rises to
$\approx580$\,ms; both flows converge to the new GEO fair share
(5\,Mbps each) within one to two seconds.
Averaged over $t \in (30, 90]$\,s, Jain's index is $J = 0.997$,
confirming near-perfect fairness.

We repeated the sweep across lead times
$\ell \in \{0, 5, 10, 15, 20, 30\}\,$s; results appear in
Table~\ref{tab:exp2} (F1~vs~BBRv3 rows).
For $\ell \leq 5$\,s and $\ell \geq 20$\,s, $J > 0.98$ in both runs.
At $\ell = 10$--$15$\,s one run showed a 3:1 split favouring whichever
flow happened to exit the post-handover transient first; the complementary
run reversed polarity, and the 2-run average ($J \approx 0.86$) reflects
the residual variance of a single 90-second trial.
The conclusion is robust: \emph{BBR-SAT does not gain a systematic
bandwidth advantage over competing BBR flows at any tested lead time.}

\subsubsection{F3: steady-state GEO}

Figure~\ref{fig:exp2_f3} shows the F3 trace.
Both flows start on the 10\,Mbps GEO beam from $t = 0$.
CUBIC reaches its characteristic sawtooth steady state by $t \approx 20$\,s.
At $t = 30$\,s BBR-SAT re-applies its stored GEO BDP context
(a no-op on link parameters, but it resets \texttt{bw\_hi},
\texttt{min\_rtt}, and \texttt{inflight\_hi} to the seeded values),
causing a brief throughput dip visible in both flows as the Jain's trace
dips below 0.8.
After $t \approx 39$\,s both flows reach a stable regime:
BBR-SAT 3.4\,Mbps, CUBIC 6.2\,Mbps, $J = 0.888$.
The 64/36 split is consistent across both runs (run-to-run $\sigma < 0.2$\,Mbps).

\subsubsection{BBR-SAT vs.\ CUBIC during handover}

We also ran the F1 topology with CUBIC as the competing flow
(Table~\ref{tab:exp2}, F1~vs~CUBIC rows).
BBR-SAT obtains only 1.1--1.9\,Mbps post-handover, with $J \approx 0.71$,
while CUBIC takes the remaining $\approx7.8$\,Mbps.
Notably, this imbalance is already present on LEO \emph{before} the
handover (CUBIC $\approx 37$\,Mbps vs.\ BBR-SAT $\approx 12$\,Mbps on
the shared 50\,Mbps LEO beam), confirming that the unfairness is
inherited directly from BBRv3's known deference to CUBIC~\cite{cardwell2017}
and is not introduced by the satellite extension.
Longer lead times ($\ell = 20$\,s) marginally worsen the split
($J = 0.655$) because BBR-SAT's proactive queue drain at $\ell$\,s before
handover vacates buffer space that CUBIC immediately reclaims.

\subsubsection{Summary}

BBR-SAT preserves the fairness properties of its BBRv3 base with
respect to both competing CCA families:
it co-exists equitably with other BBR flows ($J \geq 0.986$ at
lead times outside the narrow 10--15\,s transient window), and it
neither worsens nor repairs the pre-existing BBR-vs-CUBIC imbalance
in the steady state.
Addressing the BBR/CUBIC coexistence issue is outside the scope of
this work but is a natural direction for future study.

% ── Table ─────────────────────────────────────────────────────────────────────
\begin{table}[t]
\centering
\caption{Experiment 2 Fairness Summary (post-event average, $t > 30$\,s,
         2-run mean)}
\label{tab:exp2}
\small
\begin{tabular}{llrrrr}
\toprule
Scenario & Competitor & Lead (s) & BBR-SAT & Opponent & Jain $J$ \\
         &            &          & (Mbps)  & (Mbps)   &          \\
\midrule
F1 & BBRv3 &  0 & 4.86 & 4.84 & \textbf{0.997} \\
F1 & BBRv3 &  5 & 4.62 & 5.11 & 0.991 \\
F1 & BBRv3 & 10 & 5.83 & 3.86 & 0.862$^\dagger$ \\
F1 & BBRv3 & 15 & 3.51 & 6.17 & 0.866$^\dagger$ \\
F1 & BBRv3 & 20 & 4.62 & 5.10 & 0.987 \\
F1 & BBRv3 & 30 & 4.62 & 5.06 & 0.991 \\
\midrule
F1 & CUBIC  &  0 & 1.71 & 7.81 & 0.713 \\
F1 & CUBIC  & 10 & 1.86 & 7.80 & 0.734 \\
F1 & CUBIC  & 20 & 1.14 & 8.32 & 0.655 \\
\midrule
F3 & CUBIC  &  0 & 3.43 & 6.18 & 0.888 \\
\bottomrule
\end{tabular}
\vspace{2pt}

{\footnotesize $^\dagger$ 2-run average of high-variance single runs
(individual runs: $J = 0.738$, $J = 0.986$ and $J = 0.742$, $J = 0.991$
respectively); polarity reverses between runs, indicating random
transient ordering rather than a systematic bias.}
\end{table}
